% Сборка: latexmk лекция1_вычисления_утверждения_и_доказательства.tex
\documentclass[aspectratio=169]{beamer}
\usepackage{iftex}
\ifPDFTeX
  \usepackage[utf8]{inputenc}
  \usepackage[T1,T2A]{fontenc}
\else
  \usepackage{fontspec}
\fi
\usepackage[russian]{babel}
\usepackage{color}
\usepackage{calc}
\usepackage{graphicx}
\usepackage{epstopdf}
\usepackage{hyperref}
\hypersetup{unicode,colorlinks}
\usetheme[progressbar=head,numbering=fraction,block=fill]{metropolis}
\usepackage{minted}
\ifPDFTeX
  \usepackage{dejavu}
\else
  \setsansfont{DejaVu Sans}
  \setmonofont{DejaVu Sans Mono}
\fi
%\usepackage{adjustbox}  % Позволяет сузить куски кода (или текст) ровно настолько, чтобы уместиться в слайд
\usepackage{csquotes}
\usepackage{upquote}

\usemintedstyle{solarized-light}
% Для слайдов с блоками кода используйте \begin{frame}[fragile].
\setminted[lean4]{autogobble=true,breaklines=false}
\newminted[lean]{lean4}{fontsize=\small}
\newminted[leansmall]{lean4}{fontsize=\footnotesize}
\newminted[leantiny]{lean4}{fontsize=\scriptsize}
\newmintinline[leaninline]{lean4}{breaklines=true}

% Вывод #eval, сообщения Lean и состояния целей.
\newminted[leanoutput]{text}{autogobble=true,fontsize=\small}
\newminted[leanoutputsmall]{text}{autogobble=true,fontsize=\footnotesize}
\newminted[leanoutputtiny]{text}{autogobble=true,fontsize=\scriptsize}
\newmintinline[leanoutputinline]{text}{breaklines=true}

% Частые символы Lean в листингах при сборке pdfLaTeX.
\ifPDFTeX
  \usepackage{amssymb}
  \DeclareUnicodeCharacter{2192}{\ensuremath{\to}}
  \DeclareUnicodeCharacter{2190}{\ensuremath{\leftarrow}}
  \DeclareUnicodeCharacter{21D2}{\ensuremath{\Rightarrow}}
  \DeclareUnicodeCharacter{2194}{\ensuremath{\leftrightarrow}}
  \DeclareUnicodeCharacter{2200}{\ensuremath{\forall}}
  \DeclareUnicodeCharacter{2203}{\ensuremath{\exists}}
  \DeclareUnicodeCharacter{2227}{\ensuremath{\land}}
  \DeclareUnicodeCharacter{2228}{\ensuremath{\lor}}
  \DeclareUnicodeCharacter{00AC}{\ensuremath{\neg}}
  \DeclareUnicodeCharacter{2260}{\ensuremath{\ne}}
  \DeclareUnicodeCharacter{2264}{\ensuremath{\le}}
  \DeclareUnicodeCharacter{2265}{\ensuremath{\ge}}
  \DeclareUnicodeCharacter{2208}{\ensuremath{\in}}
  \DeclareUnicodeCharacter{22A2}{\ensuremath{\vdash}}
  \DeclareUnicodeCharacter{2115}{\ensuremath{\mathbb{N}}}
  \DeclareUnicodeCharacter{03B1}{\ensuremath{\alpha}}
  \DeclareUnicodeCharacter{03B2}{\ensuremath{\beta}}
  \DeclareUnicodeCharacter{27E8}{\ensuremath{\langle}}
  \DeclareUnicodeCharacter{27E9}{\ensuremath{\rangle}}
\fi

\vfuzz=20pt  % позволяет тексту дойти до номера слайда

\author{Алексей Романов}
\subtitle{Программирование и доказательства на Lean 4}
%\logo{}
\institute{МИЭТ}
\subject{Программирование и доказательства на Lean 4}
%\setbeamercovered{transparent}
%\setbeamertemplate{navigation symbols}{}

\title{Вычисления, утверждения и доказательства}
\date{Лекция 1}
\begin{document}
\maketitle
% --- Слайд 1. Установка Lean в VS Code или VSCodium
\begin{frame}[fragile]{Установка Lean в VS Code или VSCodium}
\begin{enumerate}
\item Установите \href{https://code.visualstudio.com/download}{VS Code} и \href{https://marketplace.visualstudio.com/items?itemName=leanprover.lean4}{расширение Lean 4}. Можно выбрать \href{https://vscodium.com/}{VSCodium} и установить \href{https://open-vsx.org/extension/leanprover/lean4}{то же расширение из Open VSX}.
\item Скачайте \href{https://github.com/alexeyr/miet-lean-course}{репозиторий курса} через \textit{Code → Download ZIP} и распакуйте архив (или клонируйте репозиторий). Через \textit{File → Open Folder} откройте папку с \texttt{lakefile.toml}.
\item Откройте \texttt{Лекция1.lean}. Для установки Lean выберите в меню \(\forall\) пункт \textit{Documentation → Show Setup Guide}.
\item В терминале редактора выполните команду \texttt{lake exe cache get} с аргументами, указанными в начале \texttt{Лекция1.lean}.
\end{enumerate}
Пока идёт загрузка, можно работать в \href{https://live.lean-lang.org/}{Lean Playground}.
\quad\href{https://lean-lang.org/install/}{Инструкция по установке}
\end{frame}

% --- Слайд 2. Первые команды в Lean Playground
\begin{frame}[fragile]{Первые команды в Lean Playground}
В \href{https://live.lean-lang.org/}{Lean Playground} выберите среду с Mathlib и начните с \leaninline|import Mathlib|. В файле курса \texttt{Лекция1.lean} найдите комментарий «Слайд 2».

Поставьте курсор на команду и посмотрите ответ в InfoView.

\leaninline|#eval 2 + 3 * 4| \qquad \uncover<2->{\texttt{14}}

\onslide<2->
Как изменится результат, если взять сумму в скобки?

\leaninline|#eval (2 + 3) * 4| \qquad \uncover<3->{\texttt{20}}

\onslide<3->
\leaninline|#check (2 + 3) * 4| \qquad \uncover<4->{\texttt{(2 + 3) * 4 : ℕ}}

\onslide<4->
При \leaninline|import Mathlib| Lean выводит тип \leaninline|Nat| как \leaninline|ℕ|.

При наведении на выражение виден его тип, а на имя — ещё и документация, если она есть.
\end{frame}

% --- Слайд 3. Выражения, типы и универсумы
\begin{frame}[fragile]{Выражения, типы и универсумы}
Запись \leaninline|e : T| читается «выражение \leaninline|e| имеет тип \leaninline|T|». Например, \leaninline|(2 : Nat)| обозначает число 2 типа \leaninline|Nat|.

Типы тоже являются выражениями. Какой тип у самого \leaninline|Nat|?

\leaninline|#check Nat| \qquad \uncover<2->{\texttt{Nat : Type}}

\onslide<2->
\leaninline|#check Type| \qquad \uncover<3->{\texttt{Type : Type 1}}

\onslide<3->
\textbf{Универсум} — это тип, элементами которого являются типы. Например, \leaninline|Nat| и \leaninline|Int| принадлежат универсуму \leaninline|Type|. Сам \leaninline|Type| принадлежит следующему универсуму \leaninline|Type 1|.

\end{frame}

% --- Слайд 4. Вычитание и типы операндов
\begin{frame}[fragile]{Вычитание и типы операндов}
Какой тип и какое значение получит каждое выражение?

\leaninline|(2 : Int) - 7|
\uncover<2->{\qquad \texttt{-5 : Int}}

\onslide<2->
Тип первого операнда задаёт тип второго: здесь \leaninline|7| имеет тип \leaninline|Int|.

\leaninline|2 - 7|
\uncover<3->{\qquad \texttt{0 : Nat}}

\onslide<3->
Для записей вроде \leaninline|2| и \leaninline|7| Lean по умолчанию выбирает \leaninline|Nat|, если их тип не задан контекстом. В натуральных числах вычитание ограничено снизу нулём.

Что произойдёт, если у операндов разные типы?

\leaninline|(2 : Nat) - (7 : Int)|
\uncover<4->{\qquad \texttt{-5 : Int}}

\end{frame}

% --- Слайд 5. Приведение типов
\begin{frame}[fragile]{Приведение типов}
Как Lean вычитает целое число из натурального?
\begin{lean}
#check (2 : Nat) - (7 : Int)
\end{lean}
\pause
\begin{leanoutput}
↑2 - 7 : ℤ
\end{leanoutput}
Знак \leaninline|↑| показывает, где Lean вставил \textbf{приведение типов}: он преобразовал натуральное число \leaninline|2| в целое. Вычитание даёт результат типа \leaninline|Int|, который mathlib обозначает \leaninline|ℤ|.

В выражении \leaninline|(2 : Int) - 7| литерал \leaninline|7| сразу получает тип \leaninline|Int| из контекста.

Наберите сокращение и нажмите пробел: \leaninput{N} для ℕ, \leaninput{Z} для ℤ, \leaninput{coe} для ↑.
Полный список в VS Code: \texttt{Ctrl+Shift+P} → \textit{Show Unicode Input Abbreviations}. См. также \href{https://github.com/leanprover/vscode-lean4/blob/master/lean4-unicode-input/src/abbreviations.json}{таблицу на GitHub}.
\end{frame}

% --- Слайд 6. Определение функции
\begin{frame}[fragile]{Определение функции}
\begin{lean}
def addOne (n : Nat) : Nat :=
  n + 1

#eval addOne 4
\end{lean}
\leaninline|def| вводит имя. В скобках стоят имя параметра и его тип. После скобок указан тип результата, после \leaninline|:=| записано тело.

\pause
При вызове \leaninline|addOne 4| параметр получает значение \texttt{4}, и команда выдаёт \texttt{5}.
\end{frame}

% --- Слайд 7. Применение и скобки
\begin{frame}[fragile]{Применение и скобки}
\begin{lean}
def square (n : Nat) : Nat := n * n

#eval square 3 + 1
#eval square (3 + 1)
#eval square (square 3)
\end{lean}
Применение функции связывает сильнее, чем сложение. Запись \leaninline|square 3 + 1| означает \leaninline|(square 3) + 1|.
\pause
Первое значение равно \texttt{10}.
\pause
Во второй строке сначала вычисляется сумма, поэтому результат равен \texttt{16}.

Что нужно вычислить первым в последней строке? Каким будет окончательный ответ? \pause Получим \texttt{9}, затем \texttt{81}.
\end{frame}

% --- Слайд 8. Функция как аргумент
\begin{frame}[fragile]{Функция как аргумент}
\begin{lean}
def twiceNat (f : Nat → Nat) (x : Nat) : Nat :=
  f (f x)

#check twiceNat addOne
#eval twiceNat addOne 3
\end{lean}
Тип \leaninline|Nat → Nat| описывает функцию с одним натуральным аргументом и натуральным результатом. Для ввода стрелки используйте \leaninput{to}.

\pause
После передачи \leaninline|addOne| остаётся функция типа \leaninline|Nat → Nat|. \pause Применив её к \texttt{3}, получим \texttt{5}.
\end{frame}

% --- Слайд 9. Функция без имени
\begin{frame}[fragile]{Функция без имени}
\begin{lean}
#eval twiceNat (fun n => n * 2) 3
#eval twiceNat (fun n => n + n + 1) 3
\end{lean}
В выражении \leaninline|fun n => n * 2| слева от \leaninline|=>| записан параметр, справа — результат. Тип параметра здесь известен из типа \leaninline|twiceNat|.

\pause
В первом примере промежуточное значение равно \texttt{6}, окончательное — \texttt{12}. Найдите оба значения во втором примере. \pause Получим \texttt{7}, затем \texttt{15}.
\end{frame}

% --- Слайд 10. Какой тип нужен для двукратного применения?
\begin{frame}[fragile]{Какой тип нужен для двукратного применения?}
\begin{lean}
def twice {α : Type} (f : α → α) (x : α) : α :=
  f (f x)

#eval twice addOne 10
#eval twice (fun b => !b) true
\end{lean}
\leaninline|α : Type| — произвольный тип (\leaninput{alpha} для α). Тип \leaninline|f : α → α| позволяет снова передать результат функции на вход.

Фигурные скобки объявляют неявный параметр. \pause
В первом вызове Lean подставляет \leaninline|Nat| и получает \texttt{12}.
\pause
Во втором — \leaninline|Bool| и получает \leaninline|true|.

\pause
Укажем тип по имени:\quad
\leaninline|#eval twice (α := Nat) addOne 10| \qquad \texttt{12}
\end{frame}

% --- Слайд 11. Что именно проверяет Lean?
\begin{frame}[fragile]{Что именно проверяет Lean?}
\begin{lean}
#eval twice addOne 3 = 5
\end{lean}
\begin{itemize}
\item \leaninline|#eval| проверяет равенство вычислением и выводит \leaninline|true|.
\end{itemize}
\begin{uncoverenv}<2->
\begin{lean}
#check twice addOne 3 = 5
\end{lean}
\begin{itemize}
\item \leaninline|#check| сообщает, что равенство имеет тип \leaninline|Prop|. Тип \leaninline|Prop| имеют и ложные утверждения.
\end{itemize}
\end{uncoverenv}
\begin{uncoverenv}<3->
\begin{lean}
  example                 : twice addOne 3 = 5 := by rfl
  theorem twice_addOne_3  : twice addOne 3 = 5 := by rfl
\end{lean}
\begin{itemize}
\item Обе записи требуют доказательства. \leaninline|example| анонимен, а \leaninline|theorem| даёт теореме имя. В следующих доказательствах можно сослаться на \leaninline|twice_addOne_3|.
\end{itemize}
\end{uncoverenv}
\end{frame}

% --- Слайд 12. Из чего состоит доказательство
\begin{frame}[fragile]{Из чего состоит доказательство}
\begin{lean}
example (n : Nat) : addOne n = n + 1 := by
  rfl
\end{lean}
Параметр \leaninline|n| произволен. За двоеточием после параметра стоит утверждение. Справа от \leaninline|:=| нужно предъявить его доказательство.

\leaninline|by| открывает блок тактик. Тактика \leaninline|rfl| закрывает равенство, если обе стороны совпадают с точностью до вычисления. Здесь достаточно раскрыть определение \leaninline|addOne|.
\end{frame}

% --- Слайд 13. Читаем цель в InfoView
\begin{frame}[fragile]{Читаем цель в InfoView}
\begin{onlyenv}<1>
\begin{lean}
example (n : Nat) : addOne n = n + 1 := by
  sorry
\end{lean}
\end{onlyenv}
\begin{onlyenv}<2>
\begin{lean}
example (n : Nat) : addOne n = n + 1 := by
  rfl
\end{lean}
\end{onlyenv}
Поставьте курсор перед тактикой. Над знаком \leaninline|⊢| видны данные и гипотезы, после него — цель.
\begin{leanoutput}
n : Nat
⊢ addOne n = n + 1
\end{leanoutput}
\leaninline|sorry| оставляет пропуск в доказательстве. Lean принимает объявление с предупреждением \texttt{declaration uses 'sorry'} и продолжает проверять файл.

\uncover<2->{Замените \leaninline|sorry| на \leaninline|rfl|. Раскрытие \leaninline|addOne| доказывает равенство. После этого шага цель закрыта, а предупреждение исчезает.}
\end{frame}

% --- Слайд 14. От вычисления к общему утверждению
\begin{frame}[fragile]{От вычисления к общему утверждению}
\begin{lean}
example {α : Type} (f : α → α) (x : α) :
    twice f x = f (f x) := by
  rfl
\end{lean}
Раскройте \leaninline|twice| в левой части. По обе стороны останется \leaninline|f (f x)|, даже когда значения \leaninline|f| и \leaninline|x| неизвестны.

Как выбор \leaninline|α| определяет типы \leaninline|f| и \leaninline|x|? Почему одно и то же доказательство \leaninline|rfl| подходит для любых \leaninline|α|, \leaninline|f| и \leaninline|x|?
\end{frame}

% --- Слайд 15. Разбираем неудачную попытку
\begin{frame}[fragile]{Разбираем неудачную попытку}
Замените правую часть предыдущего утверждения на \leaninline|f x|. После раскрытия определения тактике пришлось бы доказать
\begin{leanoutput}
⊢ f (f x) = f x
\end{leanoutput}
Проверьте кандидата в контрпримеры.
\begin{lean}
#eval twice addOne 3
#eval addOne 3
\end{lean}
\pause
Первое вычисление даёт \texttt{5}. \pause Второе даёт \texttt{4}. Здесь ошибка в самом утверждении. Верните исходную правую часть, прежде чем двигаться дальше.
\end{frame}

% --- Слайд 16. Когда утверждение верно, а шаг не проходит
\begin{frame}[fragile]{Когда утверждение верно, а шаг не проходит}
\begin{lean}
example (a b : Nat) : a + b = b + a := by
  rfl
\end{lean}
\pause
\leaninline|rfl| закрывает цель, когда обе стороны совпадают после вычисления. При произвольных \leaninline|a| и \leaninline|b| выражения \leaninline|a + b| и \leaninline|b + a| остаются разными.

\pause
\begin{lean}
example (a b : Nat) : a + b = b + a := by
  exact Nat.add_comm a b
\end{lean}
Теорема о коммутативности даёт \leaninline|Nat.add_comm a b| типа \leaninline|a + b = b + a|. Тактика \leaninline|exact| закрывает цель этим доказательством.
\end{frame}

% --- Слайд 17. Как переписать выражение с помощью rw?
\begin{frame}[fragile]{Как переписать выражение с помощью rw?}
\begin{leansmall}
example (a b : ℝ) :
    (a + b)^2 = a^2 + 2*a*b + b^2 := by
  rw [pow_two, add_mul, mul_add, mul_add, mul_comm b a,
      pow_two, pow_two, two_mul, add_mul,
      ← add_assoc, add_assoc (a * a)]
\end{leansmall}
Для вещественных чисел: \leaninline|pow_two (x : ℝ) : x ^ 2 = x * x|.

Поставьте курсор после первого \leaninline|pow_two|. В InfoView видна цель:
\begin{leanoutputsmall}
a b : ℝ
⊢ (a + b) * (a + b) = a ^ 2 + 2 * a * b + b ^ 2
\end{leanoutputsmall}
Дальше \leaninline|rw| применяет равенства по порядку. Стрелка \leaninline|←| меняет направление замены.

Какое переписывание превращает \(ba\) в \(ab\)?
\end{frame}

% --- Слайд 18. Как не перечислять все равенства?
\begin{frame}[fragile]{Как не перечислять все равенства?}
\begin{lean}
example (a b : ℝ) :
    (a + b)^2 = a^2 + 2*a*b + b^2 := by
  ring
\end{lean}
При раскрытии квадрата получаем \(a^2+ab+ba+b^2\).
Коммутативность даёт \(ba=ab\). Тактика \leaninline|ring| приводит обе стороны полиномиального равенства к общей нормальной форме и строит доказательство, которое проверяет ядро Lean.

\pause
Верна ли эта формула для матриц \leaninline|A B : Matrix (Fin 2) (Fin 2) ℝ|?

\leaninline|Fin 2| — тип с двумя значениями: 0 и 1. Первый \leaninline|Fin 2| задаёт индексы строк, второй — столбцов. Получается вещественная матрица \(2\times2\).
\end{frame}

% --- Слайд 19. Самостоятельно: два способа прибавить четыре
\begin{frame}[fragile]{Самостоятельно: два способа прибавить четыре}
\begin{lean}
def addTwo (n : Nat) : Nat := n + 2
\end{lean}
\begin{enumerate}
\item Вычислите \leaninline|twice addTwo 5|. Сформулируйте, чему равно \leaninline|twice addTwo n| для произвольного \leaninline|n : Nat|, и докажите это с помощью \leaninline|rfl|.
\item Для произвольного \leaninline|x : ℝ| запишите равенство \((x+2)+2=x+4\) и выберите тактику. Объясните, откуда взялась четвёрка.
\end{enumerate}
\pause
\leaninline|#eval twice addTwo 5| даёт \texttt{9}.
\end{frame}

% --- Слайд 20. Проверяем решения
\begin{frame}[fragile]{Проверяем решения}
\begin{lean}
example (n : Nat) :
    twice addTwo n = (n + 2) + 2 := by
  rfl

example (x : ℝ) : (x + 2) + 2 = x + 4 := by
  ring
\end{lean}
В первом случае раскрываются два определения: \leaninline|twice| и \leaninline|addTwo|. Во втором приводятся подобные слагаемые.
\end{frame}

% --- Слайд 21. От утверждения к доказательству
\begin{frame}{От утверждения к доказательству}
\begin{block}{Утверждение}
У каждого вещественного числа есть число больше него.
\end{block}

\begin{enumerate}
\item Какой тип в Lean имеет это утверждение?
\item Какие данные должен предоставить автор доказательства?
\end{enumerate}
\end{frame}

% --- Слайд 22. Что почитать и где попрактиковаться
\begin{frame}{Что почитать и где попрактиковаться}
\begin{itemize}
\item \href{https://lean-lang.org/theorem_proving_in_lean4/Dependent-Type-Theory/}{Theorem Proving in Lean 4, глава 2}. После лекции прочитайте разделы 2.1–2.4 о типах, функциях и вычислениях. Раздел 2.9 объясняет неявный аргумент \texttt{α} в нашем примере \texttt{twice}.
\item \href{https://adam.math.hhu.de/}{Lean Game Server}. Начните с \textit{Natural Number Game}, чтобы потренироваться доказывать свойства натуральных чисел.
\item \href{https://lean-lang.org/learn/}{Learn Lean}. Здесь собраны учебники по программированию на Lean и формализации математики.
\item \href{https://leanprover-community.github.io/}{Lean Community}. Здесь можно найти документацию mathlib и задать вопрос в чате Zulip.
\begin{itemize}
\item \href{https://leanprover-community.github.io/learn.html}{Учебные материалы сообщества}
\end{itemize}
\end{itemize}
\end{frame}
\end{document}
